\section{Results}

\subsection{Campaign Validity}

The main campaign executed \NumRunsMain{} runs over \NumContexts{} contexts with
\NumPolicies{} policies and \NumSeeds{} paired seeds. No run failed. Every run
completed every vehicle it generated (minimum completion \NumMinCompletion{}),
with \NumTeleports{} teleports in the entire campaign, so no result rests on a
vehicle the simulator removed. All \NumContextsAnalysed{} contexts satisfy the
validity and completeness rules declared in advance; none is excluded.

\subsection{Policy Complementarity}
\label{sec:complementarity}

The confirmatory criterion G1$'$ --- declared before the campaign and reported
here whichever way it came out --- \NumGOneVerdict{}:
\NumGOneNPolicies{} policies (\NumGOnePolicies) are strictly best in at least
one context with a seed-resolved margin. \textbf{P2 is never a resolved winner
anywhere in the factor space.} Resolved winner counts are P3 \NumWinResPThree{},
P4 \NumWinResPFour{} and P1 \NumWinResPOne{} (Table~\ref{tab:complementarity}).

The winner map is multi-factor but shallow. Naming the same policy everywhere is
correct in \NumWinMapConst{} of contexts; the best single factor
(\NumWinMapOneFactor) reaches \NumWinMapOne{}, and the best pair
\NumWinMapTwo{}. The structure is interpretable and matches the design: P1 wins
only where the shortest corridor is uncongested \emph{and} penetration is high
enough for adaptive policies to over-divert; P4 wins at low demand with generous
green; P3 wins every context at the highest demand level
(Figure~\ref{fig:regime}).

\begin{figure}[tbp]
\centering\includegraphics[width=\textwidth]{fig5_regime_map.pdf}
\caption{The regime map. Left: the resolved winner in the plane of
shortest-corridor saturation against guidance penetration; open circles mark
contexts where the seeds resolve no winner. Right: the share of contexts each
policy wins, by demand level.}
\label{fig:regime}
\end{figure}

\begin{table}[t]
\caption{Policy complementarity over 142 contexts. A win is
\emph{resolved} when the paired seed difference against the runner-up exceeds
twice its standard error. Pareto membership is over all three criteria.}
\label{tab:complementarity}
\centering\small
\begin{tabular}{@{}l r r r r@{}}
\toprule
& \multicolumn{2}{c}{Best on journey time} & \multicolumn{2}{c}{Non-dominated}\\
\cmidrule(lr){2-3}\cmidrule(lr){4-5}
& any margin & resolved & contexts & share\\
\midrule
P1 & 14 & 6 & 14 & 9.9\%\\
P2 & 42 & 0 & 74 & 52.1\%\\
P3 & 41 & 5 & 114 & 80.3\%\\
P4 & 45 & 5 & 133 & 93.7\%\\
\midrule
\multicolumn{5}{@{}l}{\emph{Pairwise differences in system mean journey time
(s/veh), positive means the first policy is slower}}\\
\midrule
Pair & mean & 5--95\% range & outcome-identical\\
\midrule
P1--P2 & +46.28 & [-1.4, +205.1] & 0.0\%\\
P1--P3 & +46.65 & [-1.6, +205.2] & 0.0\%\\
P1--P4 & +46.45 & [-1.4, +205.1] & 0.0\%\\
P2--P3 & +0.38 & [-0.5, +2.8] & 45.1\%\\
P2--P4 & +0.17 & [-0.6, +0.9] & 50.0\%\\
P3--P4 & -0.21 & [-2.4, +0.8] & 38.0\%\\
\bottomrule
\end{tabular}
\end{table}

\subsection{Decision Value of Adaptation}

Complementarity and decision value are different quantities, and here they point
in opposite directions.

The single-best fixed policy is \NumGlobalSBS{}, optimal in \NumSBSOptimalPct{}
of contexts. Figure~\ref{fig:performance} shows why: the distance-minimising policy
diverges with demand while the three adaptive policies track one another.
Running the single-best policy always costs \NumSBSMeanCost~s of system mean
journey time per vehicle; the retrospective best-per-context choice costs
\NumCostVBS~s. The headroom of Eq.~\ref{eq:headroom} is therefore
\begin{equation}
H \;=\; \NumHeadroom\ \text{s per vehicle},
\qquad
\frac{H}{\mathbb{E}_x[C(x,a_{\mathrm{SBS}})]}
\;=\; \frac{\NumHeadroom}{\NumSBSMeanCost} \;=\; \NumHeadroomPct .
\end{equation}

Selecting the wrong \emph{fixed} policy costs far more. Using the
distance-minimising policy throughout increases system mean journey time by
\NumFixedPenAbsPOne~s per vehicle relative to the single-best fixed policy, which
is \NumFixedPenPctAbsPOne{} of that policy's mean journey time of
\NumSBSMeanCost~s. The reactive and reliability-aware policies cost
\NumFixedPenPTwo~s (\NumFixedPenPctPTwo{}) and \NumFixedPenPFour~s
(\NumFixedPenPctPFour{}) on the same basis.

\begin{quote}
Choosing the fixed policy well is worth \NumFixedPenPctAbsPOne{} of mean journey
time. Choosing it \emph{per context}, given that the best fixed policy is
already in use, is worth \NumHeadroomPct{}.
\end{quote}

\subsection{Why the Adaptation Value Is Small}
\label{sec:asymmetry}

The explanation is not that the policies behave alike; Section~\ref{sec:complementarity} shows they do
not. It is that the best fixed policy is near-dominant in a measurable sense.

Define the \emph{downside} of the fixed policy as its mean loss in those
contexts where it is not best, and its \emph{upside} as its mean gain over the
best alternative in those contexts where it is. In \NumAsymNotBestPct{} of
contexts \NumGlobalSBS{} is not best, and there it loses \NumAsymLoss~s on
average (at most \NumAsymMaxLoss~s). In \NumAsymBestPct{} of contexts it is
best, and there it wins by \NumAsymMargin~s on average (at most
\NumAsymMaxMargin~s). Its upside is \textbf{\NumAsymRatio$\times$ its
downside}.

A portfolio containing such a member has a large ranking structure and a small
decision structure: any selector moving away from it risks a large loss to chase
a small gain. This is an empirical characterisation of this portfolio in this
environment, not a general theorem --- but it is measurable in any portfolio,
before a selector is built.

\begin{figure}[tbp]
\centering\includegraphics[width=\textwidth]{fig4_performance_pareto.pdf}
\caption{System mean journey time against demand at each guidance penetration
level, averaged over the remaining factors. The distance-minimising policy
diverges once the shortest corridor saturates, while the three adaptive policies
track one another closely --- the visual form of the small decision value
quantified in Section~\ref{sec:asymmetry}.}
\label{fig:performance}
\end{figure}

\subsection{Noise and Resolution of the Adaptation Effect}

Among the \NumResolvNContexts{} contexts where the fixed policy is not best, the
mean headroom is \NumResolvHeadroomNZ~s against a mean two-standard-error band
of \NumResolvNoiseNZ~s on the same paired comparison, and the headroom exceeds
its own noise in \NumResolvExceeds{} of them (\NumResolvExceedsPct{}). Resolving
the campaign-wide mean effect of \NumResolvHeadroom~s would require roughly
\NumResolvSeedsNeeded{} seeds per context, about \NumResolvRunsNeeded{} runs.
We report this rather than presenting a near-noise-floor improvement as an
achievement.

\subsection{Mechanistic Rule}

Fitted on all contexts for exposition, the depth-2 rule splits first on
$\NumMechFeature = p\,D/\mathrm{cap}_C$, the share of the shortest corridor's
capacity the guided cohort alone would consume, at a threshold bracketed between
\NumMechBracketLo{} and \NumMechBracketHi{}. The rule is physically readable:
below that share the decision turns on the effective green ratio; above it, on
whether the arterial is already oversaturated (Figure~\ref{fig:regions}).

The rule identifies the best policy in \NumMechTopOne{} of contexts, more often
than always naming the fixed policy. Its mean decision regret is nevertheless
\NumMechRegret~s against \NumMechFixedRegret~s for the fixed policy.
\textbf{It is more accurate and more costly.} Under leave-one-context-out the
same pattern holds: B1 reaches \NumBOneTopOne{} top-1 accuracy against B0's
\NumBZeroTopOne{}, and \NumBOneRegret~s of regret against \NumBZeroRegret~s.

This is the predict-then-optimise distinction~\cite{elmachtoub2022spo} in a
traffic setting. A rule trained to name the winner is rewarded for being right
where the margin is negligible and punished only mildly for being wrong where it
is large --- the wrong trade under an \NumAsymRatio$\times$ asymmetry.

\begin{figure}[tbp]
\centering\includegraphics[width=\textwidth]{fig6_selection_regions.pdf}
\caption{Selection regions in the mechanistic feature space: the true best
policy, the depth-2 mechanistic rule, and the advantage-based selector, over
controlled-flow saturation and guidance penetration.}
\label{fig:regions}
\end{figure}

\subsection{Advantage-Based Selector}

Predicting advantage (Eq.~\ref{eq:advantage}) rather than identity reverses the
failure. B4 reaches \NumBFourRegret~s of mean regret, closing
\NumBFourGapClosed{} of the gap between the fixed policy and the retrospective
best, while being \emph{less} accurate at naming the winner (\NumBFourTopOne{})
than the logistic baseline (\NumBThreeTopOne{}, \NumBThreeRegret~s). It also has
the lowest maximum regret of any selector, \NumBFourMax~s against
\NumBZeroMax~s (Table~\ref{tab:ladder}, Figure~\ref{fig:ladder}).

The primary comparison declared in advance was B4 against B1, and B4 wins it
decisively: \NumBFourRegret~s against \NumBOneRegret~s. Beating B0 alone would
not have been evidence of much, since B0 is near-dominant by
Section~\ref{sec:asymmetry}.

The magnitude is what matters operationally. B4 reduces mean regret by
\NumBFourVsBZero~s per vehicle relative to the fixed policy --- 
\NumBFourVsBZeroPct{} of a \NumSBSMeanCost~s mean journey time. The selector is
real and reproducible, and below the threshold at which an operator would act on
it.

\begin{table}[t]
\caption{Decision quality under leave-one-context-out evaluation. Regret is in
seconds of system mean journey time per vehicle. ``Within noise'' is the share
of contexts whose regret falls below that context's own seed noise floor.}
\label{tab:ladder}
\centering\small
\begin{tabular}{@{}l r r r r r r@{}}
\toprule
& \multicolumn{4}{c}{Decision regret (s/veh)} & within & gap closed\\
\cmidrule(lr){2-5}
& mean & median & max & CVaR$_{90}$ & noise & vs.\ SBS\\
\midrule
B0 ~~best fixed policy (SBS) & 0.45 & 0.00 & 17.2 & 3.7 & 99\% & ---\\
B1 ~~mechanistic rule & 0.77 & 0.00 & 17.2 & 6.5 & 96\% & -0.714\\
B2 ~~decision tree, depth 3 & 0.34 & 0.00 & 17.2 & 3.0 & 99\% & 0.238\\
B3 ~~multinomial logistic & 0.28 & 0.00 & 17.2 & 2.5 & 100\% & 0.373\\
B4 ~~GBDT advantage selector & 0.20 & 0.00 & 7.0 & 1.7 & 100\% & 0.558\\
B5 ~~selective GBDT with abstention & 0.45 & 0.00 & 17.2 & 3.7 & 99\% & 0.000\\
~~~~~~VBS (retrospective, not deployable) & 0.00 & 0.00 & 0.0 & 0.0 & 100\% & 1.000\\
\bottomrule
\end{tabular}
\end{table}

\begin{figure}[tbp]
\centering\includegraphics[width=\textwidth]{fig7_decision_ladder.pdf}
\caption{Decision quality under leave-one-context-out evaluation. Left: mean
regret for each selector; the retrospective best is hatched because it is not
deployable. Right: the share of contexts in which each selector is within the
seed noise floor, and exactly optimal.}
\label{fig:ladder}
\end{figure}

\subsection{Selective Prediction and Distribution Shift}

B5 abstains in \NumAbstainRate{} of in-distribution contexts and returns exactly
the fixed policy's regret, \NumBFiveRegret~s. The mechanism is explicit: the
conformal interval on the predicted advantage is wider than the advantage, so
the requirement that its lower bound exceed zero is never met. A calibrated gate
refuses to act on an effect smaller than its own uncertainty.

Across the \NumOODNTotal{} shift splits (Table~\ref{tab:ood},
Figure~\ref{fig:ood}), the unguarded selector B4 does not behave uniformly. Using
a tolerance of \NumOODTolerance~s per vehicle, it \textbf{helps} on
\NumOODNHelps{} splits, is \textbf{indistinguishable} on \NumOODNNeutral{}, and
\textbf{harms} on \NumOODNHarms{}:
\begin{itemize}
\item \emph{helps}: the interpolation control O3 (\NumOODThreeMidBZero{} to
\NumOODThreeMidBFour~s) and, notably, the domain shift O8
(\NumOODEightNorthBZero{} to \NumOODEightNorthBFour~s);
\item \emph{indistinguishable}: unseen low demand O2, unseen lag O4, unseen
penetration O5;
\item \emph{harms}: unseen high demand O1 (\NumOODOneHighBZero{} to
\NumOODOneHighBFour~s), unseen green ratio O6, and worst of all the concept
shift O7 (\NumOODSevenIncBZero{} to \NumOODSevenIncBFour~s).
\end{itemize}
\textbf{On every one of the \NumOODNHarms{} harmful splits, B5 recovered the
fixed policy's regret exactly.} The selective rule therefore prevented the
observed harms; it does not guarantee protection against harms not observed
here.

The support gate flags \NumOODFourLagSupport{}, \NumOODFivePenSupport{} and
\NumOODSixGreenSupport{} of held-out contexts as outside the training envelope
on the lag, penetration and green-ratio extrapolations respectively. B1 degrades
worst of all under shift, reaching \NumOODFivePenBOne~s on unseen penetration.

\paragraph{A shift the gate cannot see.} O8 is instructive precisely because the
support gate does \emph{not} fire: it flags \NumOODEightNorthSupport{} of
held-out contexts, because those contexts' descriptors are identical to
descriptors seen in training --- the incident's \emph{location} is not among the
eight. The selector happens to do well there, but it does so unguarded. A
support gate defined on a feature space cannot detect a shift that leaves that
feature space unchanged, and this is a concrete instance rather than a
theoretical caveat.

\begin{table}[t]
\caption{Distribution shift, reported separately by shift type. Regret in
seconds of system mean journey time per vehicle on the held-out regime. The
support column is the share of held-out contexts the support gate flagged as
outside the training envelope.}
\label{tab:ood}
\centering\small
\begin{tabular}{@{}l l l r r r r r r@{}}
\toprule
& Held out & Shift type & B0 & B1 & B4 & B5 & flagged & abstained\\
\midrule
O2 & demand = 1200.0 & extrapolation & 0.53 & 0.63 & 0.53 & 0.53 & 100\% & 0\%\\
O4 & lag = 300.0 & extrapolation & 0.49 & 0.42 & 0.69 & 0.49 & 100\% & 63\%\\
O5 & penetration = 0.8 & extrapolation & 1.56 & 1.01 & 0.38 & 1.56 & 100\% & 54\%\\
O6 & gc = 0.65 & extrapolation & 0.95 & 1.67 & 0.68 & 0.95 & 100\% & 86\%\\
O7 & incident = 1.0 & covariate + concept shift & 0.21 & 1.62 & 5.57 & 0.21 & 0\% & 64\%\\
\bottomrule
\end{tabular}
\end{table}

\begin{figure}[tbp]
\centering\includegraphics[width=\textwidth]{fig8_ood.pdf}
\caption{Distribution shift, reported separately by shift type. Top: mean regret
on the held-out regime for each selector. Bottom: the share of held-out contexts
flagged outside the training support, and the share on which the selective
selector abstained.}
\label{fig:ood}
\end{figure}

\subsection{Boundary Refinement}
\label{sec:boundary}

Where the main grid showed the preferred policy changing between adjacent demand
levels, three intermediate levels were run in each of the \NumBoundN{}
widest-margin cells.

\NumBoundMono{} of \NumBoundN{} refined sequences are monotone, with a single
change point, and in point estimate they narrow the switching bracket from the
\NumBoundGrid~veh/h grid spacing to a median of \NumBoundWidth~veh/h. All
\NumBoundPOneN{} shortest-to-load-balancing transitions are monotone; their
change points sit in a \NumBoundWidth~veh/h bracket whose position differs by
cell, from \NumBoundPOneEarliestLo--\NumBoundPOneEarliestHi~veh/h to
\NumBoundPOneLatestLo--\NumBoundPOneLatestHi~veh/h. That the position moves with
the cell is the point: the boundary is a property of the operating conditions,
not a single number for the network.

\textbf{None of these refined steps is seed-resolved.} The point estimates move
consistently, but no individual step between adjacent refined levels exceeds
twice its own standard error at \NumSeeds{} seeds. The remaining
\NumBoundNonMono{} sequences are non-monotone, and all involve the
reliability-aware policy. We therefore report the transition as a point-estimate
bracket and state that it is unresolved: the shortest-path to load-balancing
reversal is an ordered regime change, while the apparent reliability-aware
transitions near it are not distinguishable from noise.

\subsection{Three-Criterion Pareto Structure}
\label{sec:pareto}

Over all \NumContextsAnalysed{} contexts, \NumParetoNonSingleton{} have a
non-singleton Pareto set. Non-domination is not concentrated in one policy: P3
is non-dominated in \NumParetoPctPThree{} of contexts, P4 in
\NumParetoPctPFour{}, P2 in \NumParetoPctPTwo{} and P1 in \NumParetoPctPOne{}.

The criteria are not interchangeable. Journey time and \CO{} correlate at
\NumCorrCOneCTwo{} and rank the four policies identically in
\NumSameOrderCOneCTwo{} of contexts; journey time and stopped delay correlate at
\NumCorrCOneCThree{} and agree in only \NumSameOrderCOneCThree{}. The best
policy on \CO{} differs from the best on journey time in
\NumBestCTwoDiffersPct{} of contexts; on stopped delay it differs in
\NumBestCThreeDiffersPct{}.

The policies specialise by criterion rather than by regime.
\NumSpecCOnePolicy{} minimises journey time in \NumSpecCOneN{} contexts and
\NumSpecCTwoPolicy{} minimises \CO{} in \NumSpecCTwoN{}, but
\textbf{\NumSpecCThreePolicy{} minimises stopped delay in \NumSpecCThreeN{}
contexts (\NumSpecCThreePct{})} while minimising journey time in only
\NumSpecCOnePFour{}. The reliability-aware policy routes towards the steady,
unsignalised corridor: its vehicles spend less time stationary and more time
moving. \textbf{This is where the portfolio's complementarity lives} --- on the
criterion vector, not on the scalar cost.

\paragraph{Eco-routing remains excluded, on evidence.} \CO{} and journey time
correlate at \NumCorrCOneCTwo{} and select the same policy in
\NumSameOrderCOneCTwo{} of contexts. An objective nearly monotone in travel time
adds a label rather than a policy; the numbers are given so a reader may
disagree with the threshold rather than with an assertion.

\begin{table}[t]
\caption{Sensitivity and robustness. Upper: the two policy parameters frozen
before the experiment, the effect of changing each on that policy's own mean
journey time, and how often the identity of the best policy is unchanged.
Lower: the three declared preference profiles over normalised criteria, the
policy each selects, and how often that agrees with the journey-time choice.}
\label{tab:sensitivity}
\centering\footnotesize
\setlength{\tabcolsep}{4pt}
\begin{tabular}{@{}l l r r r r r@{}}
\toprule
\multicolumn{7}{@{}l}{\emph{Policy parameters}}\\
\midrule
Policy & Param. & Frozen & Alt. & $\Delta$ (s) & $\Delta$ (\%) & winner same\\
\midrule
P3 & \texttt{p3\_eps} & 0.2 & 0.1 & +8.10 & +1.27\% & 84\%\\
P3 & \texttt{p3\_eps} & 0.2 & 0.3 & -16.66 & -2.62\% & 78\%\\
P4 & \texttt{p4\_lambda} & 1.0 & 0.5 & -0.11 & +0.00\% & 91\%\\
P4 & \texttt{p4\_lambda} & 1.0 & 1.5 & +1.27 & +0.16\% & 98\%\\
\midrule
\multicolumn{7}{@{}l}{\emph{Preference profiles} (weights on normalised C1--C3)}\\
\midrule
Profile & Weights & P1 & P2 & P3 & P4 & agrees C1\\
\midrule
Time & (0.70, 0.15, 0.15) & 19 & 76 & 406 & 147 & 92.9\%\\
Balanced & (0.40, 0.30, 0.30) & 8 & 70 & 399 & 171 & 86.0\%\\
Environment & (0.20, 0.55, 0.25) & 7 & 66 & 425 & 150 & 80.1\%\\
\bottomrule
\end{tabular}
\end{table}

\subsection{Sensitivity}
\label{sec:sensitivity}

Across the three declared preference profiles the preferred policy changes in
\NumPrefChanges{} of contexts (Table~\ref{tab:sensitivity}). The journey-time
choice is reproduced by the time-weighted profile in \NumPrefAgreeTime{} of
contexts, the balanced profile in \NumPrefAgreeBalanced{} and the
environment-weighted profile in \NumPrefAgreeEnvironment{}. The decision is
stable under reasonable preference assumptions, and where it is not, stopped
delay is what moves it.

The two frozen policy parameters behave differently. The dispersion weight of
the reliability-aware policy is close to immaterial, changing that policy's own
mean journey time by at most \NumLambdaMaxChangePct{}; across all four parameter
changes the identity of the best policy is unchanged in \NumWinnerStableMin{} to
\NumWinnerStableMax{} of the contexts tested. The admissibility tolerance of the
load-balancing policy is not immaterial: loosening it from \NumEpsFrozen{} to
\NumEpsAlt{} improves that policy's own mean journey time by \NumEpsGainPct{}
(\NumEpsGainSec~s). The frozen value therefore \emph{understates} how dominant
the load-balancing policy is, which strengthens rather than weakens the central
conclusion. We report this rather than retrofitting the better value.
